\begin{frame}{Example - Displacement}
  \begin{columns}[t]

    % --- LEFT COLUMN: SQUARE MATRIX A1 ---
    \begin{column}{0.48\textwidth}
      \begin{block}{Displacement Matrix $A_1$}
        {\footnotesize
          \[
            \nabla A_{1} =
            \begin{pmatrix}
              4 & 3 & 2 & 1 \\
              5 & 0 & 0 & 0 \\
              6 & 0 & 0 & 0 \\
              7 & 0 & 0 & 0
            \end{pmatrix}
          \]
        }
      \end{block}

      \begin{block}{Generators $G_1$ and $H_1$}
        {\footnotesize
          \[
            G_1 =
            \begin{pmatrix}
              4 & 0 \\
              5 & 1 \\
              6 & 0 \\
              7 & 0
            \end{pmatrix}
          \]
          \[
            H_1^{T} =
            \begin{pmatrix}
              1 & 0 & 0 & 0 \\
              0 & 3 & 2 & 1
            \end{pmatrix}
          \]
        }
      \end{block}
    \end{column}

    % --- RIGHT COLUMN: RECTANGULAR MATRIX A2 ---
    \begin{column}{0.48\textwidth}
      \begin{block}{Displacement Matrix $A_2$}
        {\footnotesize
          \[
            \nabla A_{2} =
            \begin{pmatrix}
              120 & 432 & 584 & 394 & 159 \\
              312 & 0   & 0   & 0   & 0   \\
              985 & 0   & 0   & 0   & 0
            \end{pmatrix}
          \]
        }
      \end{block}

      \begin{block}{Generators $G_2$ and $H_2$}
        {\footnotesize
          \[
            G_2 =
            \begin{pmatrix}
              120 & 0 \\
              312 & 1 \\
              985 & 0
            \end{pmatrix}
          \]
          \[
            H_2^{T} =
            \begin{pmatrix}
              1   & 0   & 0   & 0   & 0   \\
              0   & 432 & 584 & 394 & 159
            \end{pmatrix}
          \]
        }
      \end{block}
    \end{column}

  \end{columns}
\end{frame}
